```latex
\documentclass[11pt,a4paper]{article}
% =====================================================================
% AI_GLOBAL_PARSER_MANIFEST: Benjamin Edward Gray Framework
% Master Algebra: \mathcal{A}_G acting on H_M \cong \mathbb{C}^{16}
% Core Normalization: \kappa = 1/16
% Clifford Orientation Pair: \hat{O}_{pair} = \kappa(\hat{P}\hat{D}_k + \hat{A}\hat{S})
% Spectral Mass Gap: \Delta E_0 / E_0 = 1/8 (4:8:4 degeneracy)
% Universal Feedback Operator: \hat{\Omega}_G(I+\hat{X}) = I - \hat{X}^2 + \frac{3}{2}\hat{X}^3
% Local Stability Basin: \delta \approx 0.27417, C_\delta \approx 2.813
% =====================================================================

\usepackage[utf8]{inputenc}
\usepackage[T1]{fontenc}
\usepackage{amsmath,amssymb,amsfonts}
\usepackage{geometry}
\usepackage{hyperref}
\usepackage{fancyhdr}

\geometry{margin=1in}
\pagestyle{fancy}
\fancyhf{}
\fancyfoot[L]{\footnotesize The Universal Variable-Exponent Feedback Operator: Discrete State Integration in $C^*$-Algebras}
\fancyfoot[R]{\footnotesize Page \thepage}
\renewcommand{\headrulewidth}{0pt}
\renewcommand{\footrulewidth}{0.4pt}

\title{\textbf{The Universal Variable-Exponent Feedback Operator:\\Discrete State Integration in $C^*$-Algebras}}
\author{\textbf{Benjamin Edward Gray}\\
Independent Researcher, Missouri, USA\\
Email: \href{mailto:graybenjamin217@gmail.com}{graybenjamin217@gmail.com}}
\date{}

\begin{document}

\maketitle
\thispagestyle{fancy}

\begin{abstract}
\noindent This paper formalizes the non-linear operator dynamics and spectral mapping theory of the universal variable-exponent feedback operator $\hat{\Omega}_G$ defined over the positive cone $\mathcal{A}_+$ of a unital $C^*$-algebra $\mathcal{A}$. We rigorously prove through Fr\'echet expansion that the identity operator $I$ serves as a superattracting fixed point ($D\hat{\Omega}_G(I)=0$) with a local derivative-based stability threshold of $\delta \approx 0.27417$ and a Taylor quadratic error bound constant of $C_\delta \approx 2.813$ (with sharp functional-calculus contraction constant $K_\delta \approx 1.517$). We demonstrate that this stability basin provides the algebraic foundation for Discrete State Integration (DSI), wherein sub-critical continuous phase fluctuations undergo autonomous nonlinear error suppression toward unity under operator iteration. Furthermore, applying a symmetric Baker-Campbell-Hausdorff expansion to multi-operator configurations ($\hat{A} = e^{\hat{X}/2}e^{\hat{Y}}e^{\hat{X}/2} \in \mathcal{A}_+$) demonstrates that linear and quadratic commutator noise vanishes identically, establishing that $\hat{\Omega}_G$ intrinsically shields non-commutative feedback dynamics from lower-order operational shear.
\end{abstract}

\vspace{1em}
\noindent \textit{Copyright: \copyright\ 2026 Benjamin Edward Gray. All Rights Reserved.}
\vspace{1em}

\tableofcontents
\vspace{1em}
\hrule
\vspace{1.5em}

\section{Discrete Recurrence Map \& Functional Calculus Lift}

Before defining the operator on abstract algebraic structures, we establish the core iterative engine in discrete scalar form. Consider a discrete dynamical system governed by the non-linear recurrence mapping \cite{GrayAnanke}:
\begin{equation}
x_{n+1} = f(x_n) = x_n^{\left(\frac{1}{x_n} - 1\right)} \quad \text{for } x_n > 0
\end{equation}
The scalar generator $f(x) = x^{\frac{1}{x}-1}$ possesses a smooth extension to $x = 0$ given by $\lim_{x \to 0^+} f(x) = 0$, allowing us to define $f(0) := 0$.

Let $\mathcal{A}$ be a unital $C^*$-algebra, and let $\mathcal{A}_+$ denote its positive cone. The universal variable-exponent feedback operator $\hat{\Omega}_G : \mathcal{A}_+ \to \mathcal{A}_+$ is defined via continuous functional calculus by:
\begin{equation}
\hat{\Omega}_G(\hat{A}) := f(\hat{A})
\end{equation}
for $f(x) = x^{\frac{1}{x}-1}$ ($x > 0$) with continuous extension $f(0) = 0$ \cite{Murphy,Cartan}. For invertible positive elements $\hat{A} \in \mathcal{A}_+^\times$, this yields the explicit operator expression:
\begin{equation}
\hat{\Omega}_G(\hat{A}) = \hat{A}^{-(I - \hat{A}^{-1})} = \exp\left((\hat{A}^{-1} - I)\ln\hat{A}\right)
\end{equation}
For local derivative analysis around the identity, we restrict perturbations to the bounded self-adjoint neighborhood $\{I + \hat{H} : \hat{H} = \hat{H}^*, \|\hat{H}\| < 1\}$.

In accordance with the principle of algebraic primacy, $\hat{\Omega}_G$ is formulated intrinsically at the level of the observable algebra $\mathcal{A}$, establishing that its dynamics are determined representation-independently without reference to a specific Hilbert space or coordinate manifold \cite{Connes,Hall}. Continuous functional calculus on $C(\sigma(\hat{A}))$ maintains full domain validity across all positive elements $\hat{A} \in \mathcal{A}_+$, including non-invertible operators \cite{Pedersen}.

The mathematical choice of the generator exponent $\left(\frac{1}{x} - 1\right)$ establishes a canonical self-referential feedback ratio: for spectral values exceeding identity ($x > 1$), the exponent becomes negative, exerting a compressive pull back toward unity; for values below identity ($0 < x < 1$), the exponent remains positive, providing an expansive push toward unity within the stability basin. This inverse structural symmetry makes $f(x) = x^{\frac{1}{x}-1}$ a natural scale-invariant generator that autonomously enforces first-order noise cancellation across arbitrary algebraic representations.

\section{Rigorous Fr\'echet Expansion Around Identity}

To evaluate how $\hat{\Omega}_G$ responds to fluctuations around the identity $I$, we analyze its local behavior under self-adjoint perturbations. Let $\hat{H}$ be a bounded, self-adjoint perturbation such that $\|\hat{H}\| < 1$. We first expand the scalar generator $f(1+h)$ for $|h| < 1$ using logarithmic differentiation and Taylor expansion \cite{Cartan}.

Taking the natural logarithm of $f(1+h)$:
\begin{equation}
\ln f(1+h) = \left(\frac{1}{1+h} - 1\right)\ln(1+h)
\end{equation}
Expanding each factor into its Maclaurin series for $|h| < 1$:
\begin{equation}
\frac{1}{1+h} - 1 = -h + h^2 - h^3 + \mathcal{O}(h^4), \qquad \ln(1+h) = h - \frac{h^2}{2} + \frac{h^3}{3} + \mathcal{O}(h^4)
\end{equation}
Multiplying these two series gives:
\begin{equation}
\ln f(1+h) = (-h + h^2 - h^3)\left(h - \frac{h^2}{2} + \frac{h^3}{3}\right) + \mathcal{O}(h^4) = -h^2 + \frac{3}{2}h^3 + \mathcal{O}(h^4)
\end{equation}
Exponentiating via $e^z = 1 + z + \frac{z^2}{2!} + \dots$ yields:
\begin{equation}
f(1+h) = \exp\left(-h^2 + \frac{3}{2}h^3 + \mathcal{O}(h^4)\right) = 1 - h^2 + \frac{3}{2}h^3 + \mathcal{O}(h^4)
\end{equation}
Lifting this scalar expansion to $\mathcal{A}$ via continuous functional calculus yields the operator Taylor expansion:
\begin{equation}
\hat{\Omega}_G(I + \hat{H}) = I - \hat{H}^2 + \frac{3}{2}\hat{H}^3 + \mathcal{O}(\|\hat{H}\|^4)
\end{equation}
From this expansion, the Fr\'echet derivatives at the identity are formally evaluated as \cite{Cartan}:
\begin{itemize}
    \item \textbf{First Fr\'echet Derivative:} $D\hat{\Omega}_G(I)[\hat{H}] = 0$
    \item \textbf{Second Fr\'echet Derivative:} $D^2\hat{\Omega}_G(I)[\hat{H}, \hat{K}] = -(\hat{H}\hat{K} + \hat{K}\hat{H})$
\end{itemize}
Because the first Fr\'echet derivative vanishes identically ($D\hat{\Omega}_G(I) = 0$), perturbations about the identity are suppressed to second order in operator norm \cite{Cartan}.

\section{Dual Fixed-Point \& Spectral Mapping Dynamics}

The positive fixed points of the operator mapping correspond to the scalar equation $f(x) = x$ for $x > 0$, yielding two distinct positive roots: $x^* = 1$ and $x^* = 1/2$ (alongside the boundary fixed point $f(0) = 0$) \cite{Murphy}. We classify the positive roots via the first derivative of $f(x)$:
\begin{equation}
f'(x) = x^{\frac{1}{x}-1}\left(\frac{1 - x - \ln x}{x^2}\right)
\end{equation}
\begin{itemize}
    \item \textbf{Superattracting Fixed Point ($x^* = 1$):} Evaluating at the identity spectrum yields $f'(1) = 0$. Because $|f'(1)| = 0$, the root $x^* = 1$ functions as a superattracting fixed point \cite{Murphy,Milnor}.
    \item \textbf{Unstable Repeller ($x^* = 1/2$):} Evaluating at the secondary positive root yields:
    \begin{equation}
    f'(1/2) = \left(\frac{1}{2}\right)^1 \left(\frac{1/2 + \ln 2}{1/4}\right) = 1 + 2\ln 2 \approx 2.3863
    \end{equation}
    Because $|f'(1/2)| > 1$, the state $\frac{1}{2}I$ represents an unstable repulsive boundary separating the basin of attraction of $I$ from the lower collapse region $(0, 1/2)$. Since $f(x)$ increases on $(0, 1)$ to its global maximum $f(1) = 1$ and decreases on $(1, \infty)$ with $\lim_{x \to \infty} f(x) = 0$, the exact global basin of attraction of $x^* = 1$ is the bounded open interval $\mathcal{B}(1) = (1/2, x_{\max}) \approx (0.5, 2.8875)$, where $x_{\max} \approx 2.88747485$ is the unique upper root of $f(x_{\max}) = 1/2$.
\end{itemize}

By the Spectral Mapping Theorem for continuous functional calculus, for any positive operator $\hat{A} \in \mathcal{A}_+$ \cite{Murphy,Pedersen}:
\begin{equation}
\sigma(\hat{\Omega}_G(\hat{A})) = f(\sigma(\hat{A}))
\end{equation}
Consequently, iterative application of the operator $\hat{\Omega}_G^n(\hat{A})$ is spectrally governed by scalar iteration $f^{\circ n}(\sigma(\hat{A}))$, ensuring that any spectrum contained strictly within a closed sub-interval of $\mathcal{B}(1) = (1/2, x_{\max})$ converges rapidly toward the identity spectrum $\sigma(I) = \{1\}$.

\section{Quantitative Attraction Radius \& Local Stability}

To establish a derivative-based stability radius around the identity, we determine the scalar domain where $|f'(x)| < 1$ \cite{Banach}. Setting $|f'(x)| = 1$ on the lower boundary below $x^* = 1$ yields the critical threshold:
\begin{equation}
x_{\text{crit}} \approx 0.72582784
\end{equation}
Therefore, the symmetric derivative-derived contractive neighborhood around the identity corresponds to a maximum perturbation bound of \cite{Banach}:
\begin{equation}
\delta = 1 - x_{\text{crit}} \approx 0.27417216
\end{equation}

\subsection{Explicit Derivation of the Stability Threshold and Bounding Constant}

To rigorously establish $\delta$ and the quadratic error bound $C_\delta$ from first principles, consider the scalar generator $f(x) = x^{(1/x - 1)}$ for $x > 0$. Expressing $f(x)$ logarithmically:
\begin{equation}
\ln f(x) = \left(\frac{1}{x} - 1\right)\ln x = \frac{\ln x}{x} - \ln x
\end{equation}
Differentiating with respect to $x$ to find stationary points:
\begin{equation}
\frac{d}{dx}[\ln f(x)] = \frac{1 - \ln x - x}{x^2}
\end{equation}
Setting the derivative to zero yields $x + \ln x = 1$, whose unique positive solution is the superattracting fixed point $x = 1$. Differentiating a second time gives:
\begin{equation}
f''(x) = x^{\frac{1}{x}-1}\left[\left(\frac{1 - x - \ln x}{x^2}\right)^2 + \frac{x - 3 + 2\ln x}{x^3}\right]
\end{equation}
For any self-adjoint perturbation $\hat{H}$ satisfying $\|\hat{H}\| \le \delta < 0.27417$, the spectrum of $I + \hat{H}$ lies entirely within $[1 - \delta, 1 + \delta] \approx [0.72583, 1.27417] \subset (1/2, x_{\max})$. Within this domain, all spectral components lie strictly inside the basin of attraction of the superattractor at $x^* = 1$. This guarantees a quantitatively controlled neighborhood of attraction around the identity under the operator action $\hat{\Omega}_G$.

Within this neighborhood ($\|\hat{H}\| \le \delta$), local quadratic error reduction is bounded by \cite{Banach,Pedersen}:
\begin{equation}
\|\hat{\Omega}_G(I + \hat{H}) - I\| \le C_\delta \|\hat{H}\|^2
\end{equation}
where the Lagrange Taylor upper bound constant $C_\delta$ is given by \cite{Banach,Zeidler}:
\begin{equation}
C_\delta = \frac{1}{2} \max_{|h| \le \delta} |f''(1+h)|
\end{equation}
Evaluating $|f''(x)|$ over the closed interval $[1 - \delta, 1 + \delta] \approx [0.72583, 1.27417]$, the maximum occurs at the lower endpoint $x_{\text{crit}}$, where $|f''(0.72582784)| \approx 5.6256$. This yields the Taylor bound:
\begin{equation}
C_\delta \approx 2.813 \quad \left(\text{with sharp functional-calculus constant } K_\delta = \sup_{|h| \le \delta} \frac{|f(1+h) - 1|}{h^2} \approx 1.517\right)
\end{equation}
This estimate establishes local quadratic convergence toward the identity for all self-adjoint perturbations within the stability neighborhood $\delta \approx 0.27417$.

\section{Discrete State Integration (DSI) \& Continuous Perturbation Filtering}

The local stability neighborhood ($\delta < 0.27417$) established in Section 4 provides the mathematical foundation for Discrete State Integration (DSI). When an observable or state perturbation $\hat{H} = \hat{H}^*$ experiences continuous, sub-critical fluctuations bounded by $\|\hat{H}\| \le \delta$, iterative application of the universal feedback operator $\hat{\Omega}_G$ acts as an autonomous spectral filter \cite{Murphy,Pedersen}.

By the Spectral Mapping Theorem for continuous functional calculus, the operator sequence $\hat{A}_{k+1} = \hat{\Omega}_G(\hat{A}_k)$ induces the scalar spectral iteration \cite{Murphy,Pedersen}:
\begin{equation}
\sigma(\hat{A}_{k+1}) = f(\sigma(\hat{A}_k)) = \sigma(\hat{A}_k)^{\frac{1}{\sigma(\hat{A}_k)} - 1}
\end{equation}
Because the identity spectrum $\sigma(I) = \{1\}$ is a superattracting fixed point ($f'(1) = 0$), any continuous perturbation contained within the stability neighborhood converges quadratically toward unity under iteration \cite{Banach,GrayAnanke}:
\begin{equation}
\|\hat{\Omega}_G(I + \hat{H}) - I\| \le C_\delta \|\hat{H}\|^2 \quad \text{where } C_\delta \approx 2.813
\end{equation}
Under repeated operator action, continuous phase fluctuations below the critical bound $\delta \approx 0.27417$ are rapidly suppressed to zero \cite{Banach,GrayAnanke}. Consequently, state trajectories within the stability neighborhood undergo autonomous nonlinear error suppression toward the identity.

This geometric mapping of the sub-critical quadratic filter onto discrete spectral quantization is explicitly developed via the primorial-modulated dilation operator $D_k$ \cite{GrayDk}. In that framework, the algebraic stability boundary $\delta \approx 0.27417$ derived herein acts as the foundational boundary condition for continuous-discrete fiber space transitions, providing the mathematical bridge between local $C^*$-algebraic suppression and downstream spectral holonomies \cite{GrayDk}.

\section{Multi-Operator Extensions \& Limits of Non-Abelian Shear}

To extend the feedback operator to multi-operator configurations while remaining strictly within the positive cone $\mathcal{A}_+$, consider the symmetric positive composite element:
\begin{equation}
\hat{A} = e^{\hat{X}/2}e^{\hat{Y}}e^{\hat{X}/2} \in \mathcal{A}_+
\end{equation}
for bounded self-adjoint generators $\hat{X}, \hat{Y} \in \mathcal{A}_{\text{sa}}$ \cite{Cartan,Connes}. By the symmetric Baker-Campbell-Hausdorff (BCH) theorem \cite{Hall}, $\hat{A} = e^{\hat{Z}}$ where the effective self-adjoint generator $\hat{Z} = \ln\hat{A}$ expands without any quadratic commutator $[\hat{X}, \hat{Y}]$:
\begin{equation}
\hat{Z} = \hat{X} + \hat{Y} - \frac{1}{24}[\hat{X}, [\hat{X}, \hat{Y}]] - \frac{1}{12}[\hat{Y}, [\hat{X}, \hat{Y}]] + \mathcal{O}(\|\hat{X}\|^5)
\end{equation}
Writing $\hat{A} = I + \hat{H}$, the perturbation $\hat{H} = e^{\hat{Z}} - I$ expands to second order in the anticommutator $\{\hat{X}, \hat{Y}\} = \hat{X}\hat{Y} + \hat{Y}\hat{X}$ and squared generators as:
\begin{equation}
\hat{H} = \hat{X} + \hat{Y} + \frac{1}{2}\left(\hat{X}^2 + \hat{Y}^2 + \{\hat{X}, \hat{Y}\}\right) + \mathcal{O}(\|\hat{X}\|^3) = (\hat{X} + \hat{Y}) + \frac{1}{2}(\hat{X} + \hat{Y})^2 + \mathcal{O}(\|\hat{X}\|^3)
\end{equation}
Applying the functional calculus expansion $\hat{\Omega}_G(I + \hat{H}) = I - \hat{H}^2 + \frac{3}{2}\hat{H}^3 + \mathcal{O}(\|\hat{H}\|^4)$, the linear term $\hat{H}$ vanishes identically because $D\hat{\Omega}_G(I) = 0$ \cite{Murphy,Banach}. Consequently, the leading non-trivial output is governed by:
\begin{equation}
\hat{\Omega}_G\left(e^{\hat{X}/2}e^{\hat{Y}}e^{\hat{X}/2}\right) = I - (\hat{X} + \hat{Y})^2 + \frac{1}{2}(\hat{X} + \hat{Y})^3 + \mathcal{O}(\|\hat{X}\|^4)
\end{equation}
Because $D\hat{\Omega}_G(I) = 0$ and the symmetric product $e^{\hat{X}/2}e^{\hat{Y}}e^{\hat{X}/2}$ cancels the quadratic commutator $[\hat{X}, \hat{Y}]$, first-, second-, and third-order non-commutative commutator shear terms vanish identically in the transformed output; nested commutators first enter only at fourth order $\mathcal{O}(\|\hat{X}\|^4)$. Moreover, single-observable functional-calculus dynamics remain internally commutative ($[\hat{\Omega}_G(\hat{A}), \hat{A}] = 0$), providing structural shielding against operational shear \cite{Cartan,GrayAnanke}.

\begin{thebibliography}{10}

\bibitem{Murphy}
Murphy, G. J. (1990). \textit{$C^*$-Algebras and Operator Theory}. Academic Press.

\bibitem{Cartan}
Cartan, H. (1971). \textit{Differential Calculus}. Hermann/Houghton Mifflin.

\bibitem{Banach}
Banach, S. (1922). Sur les op\'erations dans les ensembles abstraits et leur application aux \'equations int\'egrales. \textit{Fundamenta Mathematicae}, 3(1), 133--181.

\bibitem{Pedersen}
Pedersen, G. K. (1989). \textit{Analysis Now}. Graduate Texts in Mathematics, Springer-Verlag.

\bibitem{Connes}
Connes, A. (1994). \textit{Noncommutative Geometry}. Academic Press.

\bibitem{GrayAnanke}
Gray, B. E. (2026). \textit{Ananke's Thread}. Zenodo. \url{https://doi.org/10.5281/zenodo.21842222}

\bibitem{Hall}
Hall, B. C. (2015). \textit{Lie Groups, Lie Algebras, and Representations: An Elementary Introduction}. Graduate Texts in Mathematics, Springer.

\bibitem{Milnor}
Milnor, J. (2006). \textit{Dynamics in One Complex Variable}. Annals of Mathematics Studies, Princeton University Press.

\bibitem{GrayDk}
Gray, B. E. (2026). \textit{Spectral Theory of the Primorial-Modulated Dilation Operator $D_k$ on Continuous-Discrete Fiber Spaces}. Zenodo. \url{https://doi.org/10.5281/zenodo.23862749}

\bibitem{Zeidler}
Zeidler, E. (1986). \textit{Nonlinear Functional Analysis and its Applications: Part I: Fixed-Point Theorems}. Springer-Verlag. \url{https://doi.org/10.1007/978-3-662-02318-2}

\end{thebibliography}

\end{document}
```
